\documentclass[10pt]{article}
\usepackage{template}
\usepackage{amsmath,amssymb,microtype}
\setlength{\emergencystretch}{2em}
\SetKeys{date=2026-10-18,vol=005,title={每日有趣新闻｜科学与工程},author={ChatGPT 编辑部},note={\small 2026 年 10 月 18 日晨刊｜北京时间 09:50 截稿。重点报道 10 月 17 日的新进展；对于较早发表、昨日重新受到关注的论文，明确区分论文日期与报道日期。}}
\renewcommand{\newsitem}[3]{\par\medskip{\large\bfseries\color{newsblue}#1\par}{\footnotesize\color{newsgray}#2\par}\smallskip\color{black}#3\par\medskip{\color{newsgray}\rule{\linewidth}{0.25pt}}\par}
\newcommand{\newslead}[3]{\par\medskip{\LARGE\bfseries\color{newsblue}#1\par}\vspace{1mm}{\small\color{newsgray}#2\par}\smallskip\color{black}#3\par\medskip{\color{newsgray}\rule{\linewidth}{0.45pt}}\par}
\begin{document}
\begin{documents}

\newslead{固态硬盘说自己已经就绪，内核却不能信：Linux 7.3 最后阶段的 PCIe 修补}{头条｜Linux 内核、PCIe、虚拟化｜10 月 17 日报道；补丁讨论始于 10 月 11 日}{
Linux 7.3 正式版发布前夕，一项针对三星 990 PRO 与 990 EVO Plus 固态硬盘的兼容性修补进入内核。乍看只是设备表中增加了一个例外，但它揭示了操作系统与硬件之间很重要的边界：驱动程序通常依赖设备报告的能力来决定如何操作；一旦设备宣称的能力与固件实际行为不一致，原本合理的优化就可能演变为功能故障。10 月 17 日，Phoronix 报道这一修补已随 PCI 子系统修正进入 Linux 7.3 的发布准备阶段。\cite{PCIreport}

\textbf{问题发生在重置之后。}PCI Express 的功能级重置（Function Level Reset，FLR）允许软件在不重启整个系统的情况下重置一个 PCIe 功能。对于把实体设备交给虚拟机直接使用的 VFIO 直通场景，这是一项关键机制：虚拟机退出、设备被释放，再启动另一台虚拟机时，宿主机需要把设备恢复到可预测的状态。重置不等于立刻可用，硬件可能需要一段时间完成内部初始化。PCIe 规范提供了名为 Immediate Readiness 的能力指示，用来告知软件是否可以省略通常的等待。

三星部分 990 PRO 固件报告支持 Immediate Readiness，但在 FLR 后最初约七毫秒，配置空间的读取结果可能仍表现为全一。内核于是面临一个矛盾：设备自己说可以立即响应，实际却尚未准备好。内核此前对恢复基址寄存器（BAR）过程做过优化，减少了多余的重试；两者叠加之后，重置后的配置恢复可能丢失。提交者 Alexander Gruhlke 报告，受影响配置在第二次启动直通虚拟机时可能失败，严重时需要重新启动宿主机才能重新使用设备。这里的“第二次”很关键，因为它把问题定位到设备重置与恢复流程，而不是首次发现硬盘或日常读写。\cite{PCIpatch,PCIdiscuss}

\textbf{为何不简单要求更新固件？}补丁讨论列出，990 PRO 固件 8B2QJXD7 可以在重置后立即响应，而旧版本 1B2QJXD7 存在上述延迟；然而 PCI 设备修正代码在相应阶段不一定能获取足够的固件版本信息。维护者必须权衡：仅提醒用户升级固件，还是为整个设备型号禁用这项能力。后一种做法可能使正常固件也承担额外等待，却能避免已知的恢复失败。公开讨论同时提到 990 EVO Plus 的类似问题；关于不同型号和固件的适用范围，应以最终合入的设备匹配表为准，不能把一个型号的测试结果推广为所有三星 NVMe 产品的通病。\cite{PCIdiscuss,PCIreport}

这一修补的思路是针对相关设备忽略错误的 Immediate Readiness 宣称，让内核按照需要等待设备恢复。它不是提高顺序读写速度的优化，也不表示这些硬盘会在普通 Windows 使用中突然损坏数据。受影响的核心场景是 Linux 主机上的 PCIe 重置，尤其是 VFIO 设备直通；是否实际触发还与内核版本、设备固件和使用方式有关。不能把“某一驱动路径存在兼容性问题”扩大为“整个系列 SSD 不可靠”。

\textbf{更普遍的工程教训。}接口标准规定了设备与软件应该如何协作，但软件无法单凭类型定义或规范文本保证真实硬件始终遵守约定。内核中的 quirk（针对特定硬件的例外处理）看似不够优雅，却承担着把理想化协议适配到真实产品的任务。修补之所以值得关注，还因为它说明了回归如何产生：某项优化在符合规范的设备上完全正确，却可能暴露长期潜伏的固件缺陷。验证系统级优化时，除性能基准外，重置、热插拔、错误恢复与虚拟化生命周期都应进入测试矩阵。

本刊核查了 10 月 11 日的内核邮件列表补丁与 10 月 17 日的技术报道。尚未对具体硬盘进行独立复现实验；因此本文描述的是公开报告和补丁解释，而非本刊实测。读者若正在排查 VFIO 直通故障，应先核对设备的 PCI ID、固件版本、内核版本与日志，再决定是否更新固件或等待包含修正的内核包。
}

\newslead{728 个量子比特测量三维拓扑：比规模数字更值得看的，是如何验证体态}{深度报道｜量子模拟、凝聚态物理｜10 月 17 日报道；预印本 10 月 12 日提交}{
量子计算新闻常把量子比特数量放在标题最醒目的位置。然而，10 月 17 日受到关注的一项 D-Wave 研究，最有意思的部分并不是使用了多达 728 个超导量子比特，而是研究团队尝试在可编程模拟器中同时读取拓扑体系的边界特征和体态不变量。预印本《Measuring higher-order bulk topology in a programmable quantum simulator》于 10 月 12 日提交至 arXiv，17 日有行业媒体作了进一步报道。本文将它视为近期论文的专题解读，而不误写成 17 日才发表的新实验。\cite{Topopaper,Toporeport}

\textbf{拓扑与通常的相变有什么不同？}传统凝聚态物理经常使用局域序参量区分物态。例如铁磁体中，自旋在宏观尺度上趋向同一方向，磁化强度便可以表征有序程度。拓扑绝缘体的关键性质则不能简单由某一点的局域量说明，而与量子态在整个参数空间中的几何结构有关。体系内部可以呈现能隙，边缘或表面却可能支持特殊的局域态。所谓“体边对应”，就是把体内的拓扑性质与边界上可观察的特征联系起来。

高阶拓扑绝缘体进一步改变了边界态出现的位置：二维体系可能在角点形成局域态，三维体系可能出现棱边或角点态。直观地说，普通拓扑绝缘体常讨论比体低一个维度的边界，而高阶拓扑体系允许特殊态集中在更低维的位置。这里的“高阶”并不意味着量子计算机性能等级更高，而是描述边界态的维度层级。

\textbf{为什么只看角点还不够？}实验中发现角点附近的局域信号，并不能自动证明它来自所宣称的拓扑机制。缺陷、势阱和有限尺寸效应也可能制造局域态。因此，研究者希望独立获得体内的拓扑信息，再检查它是否与边界态相符。论文使用实空间测量重构复传播子，从中推断有效能带投影算符，并进一步计算嵌套 Wilson 环路。Wilson 环路可以理解为量子态沿闭合路径演化时积累的几何信息；在高阶拓扑体系中，对相关子空间继续进行嵌套分析，能够提取与多极矩有关的不变量。

团队在一维、二维与三维有效紧束缚模型中进行实验，报告了局域态密度所呈现的边界信号，以及从体态重构得到的拓扑表征。在三维模型中，研究者给出了通过第三层嵌套 Wilson 环路提取的八极矩。与只展示边界上的一个亮点相比，体内与边界两条证据链相互印证，才是这项研究在方法论上的进展。\cite{Topopaper}

\textbf{728 个量子比特到底代表什么？}论文中的处理器是用于量子模拟的超导模拟—数字混合平台。研究者通过编程构造有效哈密顿量，观察所需的动力学和谱学特征。这与一台可以稳定执行任意量子算法、具有完整容错纠错机制的通用量子计算机不是同一回事。量子比特总数也不能直接换算成可用逻辑量子比特数，更不能凭此宣称已经实现对经典计算机的普遍超越。

预印本仍需经历同行评议和其他团队的独立验证。对于模拟器实验，关键问题包括：有效模型与目标拓扑模型的偏差有多大，控制误差如何传播到重构出的投影算符，噪声和有限尺寸会否改变不变量的判断，以及不同参数下结果是否稳定。报道没有给出足够材料让本刊独立复算这些不变量，因此我们不把预印本的结论表述为已由外部实验室确认。

从更长远的视角看，可编程量子模拟器的价值可能不仅在于“模拟出一种物态”，还在于发展可重复、可比较的测量协议。若体态重构的方法能够扩展到其他模型，它将帮助研究者把量子模拟从现象展示推进到定量检验。但这种前景仍是研究方向，而不是已经完成的产业应用。
}

\newsitem{启动日志线程也会影响多核初始化：Linux 修补罕见开机卡死}{重点报道｜Linux 内核｜10 月 17 日报道；原补丁 10 月 15 日}{
开机过程中输出日志似乎只是一个旁观操作，但 Linux 的非阻塞控制台（NBCON）实现说明，记录系统状态本身也会参与系统调度。内核开发者 John Ogness 在 10 月 15 日提交修补，拟推迟控制台打印线程的启动时间；10 月 17 日，Phoronix 报道相关修正进入 Linux 7.3 的最后阶段。补丁仅修改了初始化调用的级别，却对应一些难以稳定复现的多处理器开机卡死报告。\cite{NBpatch,NBreport}

问题集中在对称多处理器（SMP）初始化的早期阶段。内核需要依次启动其他处理器核心，并通过处理器间中断（IPI）完成同步。在这段敏感时间内，控制台打印线程如果异步启用、禁用中断，就可能与平台的启动时序产生交互。维护者观察到，一些故障报告在撤回 NBCON 控制台使用后消失；这构成了可疑线索，却不意味着每一次开机卡死都由同一原因引起。公开邮件还指出，部分受影响硬件可能存在固件或微码层面的因素。\cite{NBpatch}

修补把打印线程的启动从 \texttt{early\_initcall()} 延后到 \texttt{subsys\_initcall()}，也就是让部分架构初始化工作先完成。其工程判断是：稍晚开始异步打印，通常比让系统在启动期间失去响应更可接受。这与许多并发错误的特点相符：单个操作并没有明显错误，但两个合法操作在特定时间窗口内交错，就可能触发隐藏的时序假设。

从调试角度看，这类问题尤其棘手。添加日志会改变线程调度与中断时序，结果可能使问题更难复现；不同 CPU、固件版本、启动参数又会影响发生概率。因此，维护者需要结合回归范围、日志路径、硬件报告与重复启动测试，而不能把偶发成功当成修复证明。补丁讨论中的测试仍有待持续验证，后续稳定内核分支是否回移，也应等待相应维护者正式发布。

这则新闻提醒开发者，所谓“辅助功能”并不必然与核心行为隔离。日志、监控、性能采样甚至异常处理，都可能在并发系统中引入新的同步边界。对复杂设备或驱动程序而言，能够观察故障与不会干扰故障，常常是两件不同的事。
}

\newsitem{GNOME 风格的 Moose 更新：本地模型与云端模型之间，权限应怎样划线？}{重点报道｜Linux 桌面、Rust、AI｜10 月 17 日正式发布 v0.6.1}{
开源桌面 AI 客户端 Moose 于 10 月 17 日发布 v0.6.1。根据项目维护者的发行说明，新版增加了 Mistral AI 与 DeepSeek 连接能力，也允许用户按照受支持的 API 格式自行添加服务提供商；模型可以自动发现，也可手动填写标识符，部分模型还能单独配置图像输入及推理选项。新版同时修复了托管 Ollama 安装路径包含符号链接时的安装问题。上述信息来自项目自己的版本记录，而非仅根据第三方截图推断。\cite{Moose}

理解这次更新，还需要把它放在 10 月 7 日的 v0.6.0 基础上看。上一版本引入了云端 API 密钥的受保护存储、文件浏览器，以及只使用本机托管 Ollama 的 Local Only 模式；用户可按对话设置分享权限，消息与文件拥有不同的权限开关。项目文档还明确说明，聊天记录、草稿和文档库保存在本机，而主动发送给远程提供商的内容仍受对方数据处理政策约束。\cite{Moose06}

这意味着“支持本地模型”和“数据永远不离开设备”不能划等号。一个同时支持本地与远程提供商的客户端，必须在界面和内部权限模型中清楚地区分两个数据路径。若用户上传了源代码或实验数据，决定是否发送到云端的应当是明确的授权，而不只是一个容易忽略的模型名称。对于多模型聊天软件，文件权限、消息权限、模型选择和历史记录归属最好都能被用户检查。

Moose 使用 GTK4、libadwaita 等 GNOME 技术栈，但它是面向 GNOME 桌面体验的第三方应用，不能据此说它已成为 GNOME 核心组件。此次更新也没有证明其远程连接的隐私机制经过独立安全审计。评估这类工具时，应同时看功能完整度、密钥存储方式、实际网络请求与可审计的代码，而不是仅凭“开源”“本地优先”等宣传词。
}

\newsitem{十年前的彗星照片，今天为什么仍能成为立体科学图像？}{天文学与科学图像｜NASA APOD 10 月 17 日刊登；原始影像摄于 2015 年}{
NASA 的每日天文图片栏目在 10 月 17 日刊登了一幅彗星 67P／丘留莫夫—格拉西缅科的红蓝立体合成图。图像来自欧洲空间局 Rosetta 探测器搭载的 OSIRIS 窄角相机：原始照片拍摄于 2015 年 7 月 25 日，距离彗核约 184 千米。两张从略有差异的视角取得的图像被组合为红蓝视差图，配合相应眼镜可以让左右眼接收不同视角的信息，从而获得立体深度感。这里 2026 年 10 月 17 日是栏目发布日，不是探测器的新观测日期。\cite{APOD67P}

67P 的形状很有辨识度：较大的彗核部分直径约四千米，较小部分约二点五千米，中间由狭窄的连接区域相连。彗星接近太阳时，表面或近表层的挥发性物质受热释放气体，并裹挟尘埃形成喷流。照片中能够看见喷流，并不意味着一张静态图就能直接测出它们的三维速度；那需要结合不同时间的影像、成像几何与物理模型。\cite{APOD67P}

立体图像的科学价值在于，视差可以帮助区分真正的地形起伏与光照阴影造成的错觉。对形状极不规则的小天体而言，地形几何会影响太阳照射角、热量分布、气体逸出方向以及表面物质的迁移。只靠一张二维图，明暗变化容易被误认为高度变化；两个视角则提供额外的几何约束。

Rosetta 任务在 2016 年 9 月以受控撞击彗星表面结束。十年后重新展示旧数据，不等于宣布发现了新彗星；它说明空间探测器留下的高质量原始影像仍可以通过不同的可视化方式帮助公众理解天体结构。本文只对 NASA 的图像说明作技术解读，未独立进行 OSIRIS 影像的三维重建。
}

\newsitem{观测预告｜猎户座流星雨即将迎来峰值}{天文简讯｜NASA 2026 年观测指南}{
NASA 的 2026 年观测指南列出，猎户座流星雨活跃期约为 10 月 2 日至 11 月 7 日，预计在 10 月 21 日至 22 日夜间达到峰值。它与哈雷彗星留下的尘埃轨道有关，流星进入地球大气层的典型速度约为每秒 66 千米。理想暗夜环境下可见流星数量通常为每小时十至二十颗，但具体观测还受到月光、天气、光污染与辐射点高度影响。此处属于未来观测预告，而非过去一天发生的新发现。\cite{Orionids}
}

\newsitem{内核维护者亲自使用 LLM 改造页缓存：二十三倍成绩与一个竞态漏洞}{技术专题｜Linux、并发算法、AI 编程｜10 月 17 日仍在讨论；代码系列始于 10 月 13 日}{
本周 Linux 内核社区还有一项值得从方法论上观察的实验：Linus Torvalds 在 10 月 13 日公开了一组针对页缓存引用计数瓶颈的补丁，说明自己在开发过程中使用了大语言模型协助。10 月 17 日，相关技术媒体仍在讨论这一系列补丁及其审查结果。原始邮件把问题追溯到多核处理器同时读取同一份缓存文件数据时的共享引用计数争用。它不是 17 日才提交的全新补丁，更不是已进入稳定内核的最终成果。\cite{LinusRCU}

Linux 将文件内容缓存为页或更大的 folio。通常，读取路径需要确保读者使用数据期间，底层内存不会被另一条路径释放。取得引用计数是常见保护方式，但当几十个 CPU 同时修改同一个 folio 的引用计数时，缓存行在不同核心间迁移，原本很轻的读取也会出现明显的同步成本。新方案试图把部分频繁访问的 folio 标记为延迟释放，在满足条件时借助 RCU（读—复制—更新）保护读取窗口，使热路径可以跳过部分引用计数操作。

RCU 的思想并不是让内存永远不释放，而是把“从数据结构中移除对象”和“真正回收其内存”分开，中间留出一个宽限期，确保旧读者已经退出受保护的访问区间。这在读多写少的数据结构中经常有优势；但若要把它应用到文件内容读取，还必须同时考虑文件截断、页替换、缺页异常、缓存失效以及内存压力。性能提升的代价可能是更复杂的对象生命周期证明。\cite{LinusRCU}

Torvalds 报告，在一项专门让大量线程反复读取同一小块缓存数据的合成测试中，六十四线程吞吐量由每秒约 1330 万次升至 3.04 亿次，约为二十二点九倍。他本人强调，这种极端构造的测试不代表真实应用整体能够加速二十三倍。对于数据库或文件服务器而言，缓存命中率、存储延迟、锁竞争、工作集大小与用户态处理都会共同决定收益，不能从一个微基准直接外推。\cite{LinusRCU}

更重要的是，公开审查随后发现了竞态条件：如果读者持有的 folio 在访问期间被替换，同时文件长度又发生变化，原有检查顺序可能使读者得到不应暴露的数据。审查者提出了额外的一致性检查，并继续讨论读取顺序和内存屏障。这个过程既证明自动化工具可以参与复杂系统的原型开发，也证明一个漂亮的性能数字并不包含正确性证明。对内核代码来说，安全性必须通过人类审查、可复现测试和严谨的并发推理共同建立。\cite{LinusRace}

这起事件还值得与前文的 SSD 和启动日志修补对照：三者都与时间窗口有关，却分处不同层次。PCIe 问题是设备重置后未及时响应，NBCON 问题是早期线程和中断交错，页缓存问题则是多个 CPU 对对象生存期的竞争。若只把这些案例归结为“某个程序员写错了一行”，便会错过系统工程真正困难的部分：不同参与者对同一时刻的状态拥有不同假设。
}

\newsitem{九十五年前的 Bethe 弦，如何在超冷原子气体中留下可测的证据？}{量子物理深读｜10 月 17 日媒体再报道；研究机构 9 月 14 日已公布}{
10 月 17 日的一篇科普报道重新介绍了因斯布鲁克大学团队对 Bethe strings（Bethe 弦）的实验研究。它的原始大学新闻稿发表于 9 月 14 日，相关论文刊于《Nature Communications》；因此这不是昨日刚完成的实验，而是一项近期成果在周末再次受到关注。\cite{Bethe,BetheReport}

1931 年，Hans Bethe 在求解一维量子多体模型时发展出一套特殊方法。某些相互作用条件下，多粒子态的数学解可以组成所谓的“弦”，表现为一组具有特定关联结构的复数准动量。它们不应被想象为真的用绳子把原子串在一起；这个名称描述的是可积模型中多体束缚态的解结构。在实验上，研究者关心的是能否用可控的原子体系实现相应的集体状态，并从可观测量识别其束缚特征。

团队先将铯原子冷却至绝对零度以上数十亿分之一开尔文的温度，再用光学势把气体分隔到大量近似一维的细管中。磁场调节原子间相互作用，使体系从排斥条件进入吸引条件。在合适的密度与相互作用强度下，原子可以形成由多个粒子共同构成的束缚态；大学说明提及包含六个或更多原子的簇。这里的束缚不是通常化学键形成分子的过程，而是特殊低维量子多体体系的相互作用结果。\cite{Bethe}

如何知道这些簇真的存在？团队采用了两种不同的释放方式。若仍把原子限制在一维管内，让它们沿管道扩展，束缚态在碰撞中可以保持相对稳定；若突然撤去限制，让体系向三维自由空间扩展，原先只在一维条件下成立的束缚会被破坏，释放出的束缚能转化为额外的运动能量。比较两种情况下的扩展能量，便能得到关于束缚态存在的间接证据。对于排斥相互作用的非束缚气体，这种额外能量差异则不应以相同方式出现。

实验与理论的符合程度，是这项工作的主要价值。但必须区分“通过动力学观测推断束缚态”与“逐个拍到每一条弦的内部结构”。现有测量不是对每个原子的位置进行连续直接成像；它提供的是受控释放实验与理论模型相互印证的证据。后续如果能进一步研究不同簇尺寸、碰撞寿命和非平衡演化，才有机会检验更多关于可积系统的精细预言。\cite{Bethe}

这则研究之所以有趣，也因为它把数学物理中相当抽象的结构转化成了实验可检验的问题。它不意味着所有一维材料都会自发形成相同束缚态，更不意味着发现了新的基本粒子。真正的进展是获得了一套可调几何维度、相互作用与粒子密度的实验平台。
}

\newsitem{一层原子的台阶，能把超导涡旋变成沿轨道运动的对象吗？}{凝聚态物理｜10 月 16 日再报道；论文 7 月 30 日发表}{
日本物质材料研究机构的团队在原子层超导体系中观察到显著的涡旋输运各向异性。研究对象是带有规则原子台阶的表面，扫描隧道显微镜帮助确认台阶和局域涡旋结构；四端电阻测量则显示，在一定磁场区间内，沿台阶方向与跨越台阶方向的有效涡旋迁移率可以相差约三个数量级。论文于 7 月 30 日发表在《Physical Review B》，10 月 16 日有新的科普报道，并非本周才获得的原始实验结果。\cite{VortexPaper,VortexReport}

超导体的“零电阻”并不表示内部完全没有磁通结构。某些条件下，磁通以量子化涡旋形式穿透材料，周围伴随超导电流。当电流或其他作用力推动涡旋运动时，涡旋运动可能引起耗散，进而影响测得的电阻。因此，理解并控制涡旋钉扎与运动方向，是超导器件研究中的重要课题。

原子台阶可以改变局部耦合环境，让涡旋沿不同方向运动时面对不同的能量障碍。研究者在中等磁场范围报告了约千倍的各向异性，并在大约 0.10 至 0.20 特斯拉的条件下识别出沿台阶方向近似无钉扎的一维涡旋流动区域。在更低温度下，实验数据还显示与量子隧穿相关的输运机制。这里的千倍并不是整个超导器件性能提升千倍，而是特定条件下的涡旋迁移率比值。\cite{VortexPaper}

这一研究可以被理解为利用表面几何来塑造涡旋的运动景观。但“让涡旋容易运动”未必总是好事：对于希望降低耗散的应用，研究者有时恰恰需要把涡旋牢牢钉住；对于希望操控涡旋或研究量子动力学的器件，定向移动则可能更有价值。未来能否制造可重复、稳定且适用于实际电路的原子台阶结构，还需要材料生长、器件集成与工作温区方面的验证。
}


\newsitem{编者说明｜本期为什么保留“论文日期”和“报道日期”两套时间？}{编辑札记｜2026-10-18}{
本期最即时的工程新闻是 10 月 17 日对 Linux 内核修补的报道，以及同日发布的 Moose 版本；量子拓扑论文在 10 月 12 日已提交预印本，Bethe 弦实验与超导涡旋研究更早发表，而彗星图像来自 2015 年。把这些材料统称为“昨天的新发现”会制造错误的时间感。因此本刊将当日进展、旧研究的最新报道与观测预告分别标注。技术新闻中的性能数字和实验指标，亦须对应其原始测试条件，不能自动外推到所有设备或实际工作负载。
}

\par\medskip{\large\bfseries\color{newsblue}参考资料与原始出处}\par
\begingroup\footnotesize\renewcommand{\refname}{}\bibliographystyle{unsrt}\bibliography{refs}\endgroup
\end{documents}
\end{document}